% !TEX program = xelatex
\documentclass[12pt,a4paper]{ctexart}

% ============================================================
% 页边距与排版
% ============================================================
\usepackage[top=2.5cm, bottom=2.5cm, left=2.5cm, right=2.5cm]{geometry}
\usepackage{setspace}
\onehalfspacing

% ============================================================
% 数学与物理
% ============================================================
\usepackage{amsmath, amssymb, amsthm}
\usepackage{physics}
\usepackage{braket}
\usepackage{bm}

% ============================================================
% 图表与颜色
% ============================================================
\usepackage{graphicx}
\usepackage{float}
\usepackage{booktabs}
\usepackage[table]{xcolor}
\usepackage{listings}
\usepackage{tcolorbox}
\tcbuselibrary{listings, breakable, skins}

% ============================================================
% 超链接与参考文献
% ============================================================
\usepackage{hyperref}
\hypersetup{
    colorlinks=true,
    linkcolor=blue,
    citecolor=blue,
    urlcolor=blue
}

% ============================================================
% 定理环境
% ============================================================
\newtheorem{theorem}{定理}[section]
\newtheorem{definition}{定义}[section]
\newtheorem{remark}{注记}[section]
\newtheorem{equationN}{公式}[section]

% ============================================================
% 代码环境设置
% ============================================================
\definecolor{codebg}{RGB}{245,245,245}
\definecolor{codeframe}{RGB}{200,200,200}
\definecolor{keyword}{RGB}{0,0,180}
\definecolor{comment}{RGB}{0,130,0}
\definecolor{string}{RGB}{180,60,0}
\definecolor{number}{RGB}{120,120,120}

\lstdefinestyle{pythonstyle}{
    language=Python,
    basicstyle=\ttfamily\small,
    backgroundcolor=\color{codebg},
    frame=single,
    rulecolor=\color{codeframe},
    breaklines=true,
    breakatwhitespace=false,
    numbers=left,
    numberstyle=\tiny\color{number},
    numbersep=8pt,
    keywordstyle=\color{keyword}\bfseries,
    commentstyle=\color{comment}\itshape,
    stringstyle=\color{string},
    showstringspaces=false,
    tabsize=4,
    captionpos=b,
    xleftmargin=0.5cm,
    xrightmargin=0.5cm,
    aboveskip=8pt,
    belowskip=8pt
}

\lstdefinestyle{bashstyle}{
    language=bash,
    basicstyle=\ttfamily\small,
    backgroundcolor=\color{codebg},
    frame=single,
    rulecolor=\color{codeframe},
    breaklines=true,
    numbers=left,
    numberstyle=\tiny\color{number},
    numbersep=8pt,
    keywordstyle=\color{keyword}\bfseries,
    commentstyle=\color{comment}\itshape,
    showstringspaces=false,
    tabsize=4,
    captionpos=b,
    xleftmargin=0.5cm,
    xrightmargin=0.5cm,
    aboveskip=8pt,
    belowskip=8pt
}

% 带标题的代码框
\newtcblisting{codeblock}[2][]{
    breakable,
    colback=codebg,
    colframe=codeframe,
    listing only,
    listing options={style=pythonstyle},
    title={\textbf{#2}},
    #1
}

\newtcblisting{bashblock}[2][]{
    breakable,
    colback=codebg,
    colframe=codeframe,
    listing only,
    listing options={style=bashstyle},
    title={\textbf{#2}},
    #1
}

% ============================================================
% 自定义命令
% ============================================================
\newcommand{\Pz}{P^{z}}
\newcommand{\tsep}{t_{\text{sep}}}
\newcommand{\QCD}{\Lambda_{\text{QCD}}}
\newcommand{\hR}{\tilde{h}_R}
\newcommand{\code}[1]{\texttt{#1}}
\newcommand{\physeq}{\ensuremath{\longleftrightarrow}}
\newcommand{\refpaper}[1]{参考文献[\ref{#1}]}

% ============================================================
% 标题页
% ============================================================
\title{\textbf{非极化胶子PDF的格点QCD计算}\\[0.3cm]
       \large huangcl代码的物理公式与程序实现对应分析}
\author{张鑫 \\ 中国科学院近代物理研究所}
\date{\today}

\begin{document}

\maketitle

\begin{abstract}
本文对 huangcl 编写的非极化胶子 Parton Distribution Function (PDF) 格点计算代码
(\texttt{examples/huangcl/code.py}) 进行系统的物理分析。
我们从大动量有效理论 (LaMET) 出发，逐层建立物理公式与程序实现之间的对应关系，
涵盖场强张量的格点离散化、非定域胶子算符的 OPE 构造、蒸馏法核子两点关联函数、
非连通三点关联函数的分解、比值方法、以及 Jackknife 重采样误差分析。
每个物理公式均附有对应的 Python 实现，使物理图像与编程图像清晰对照。

\textbf{关键词：}胶子 PDF，LaMET，格点 QCD，场强张量，非连通图，OPE，Jackknife
\end{abstract}

\newpage
\tableofcontents
\newpage

% ============================================================
% 第一章：引言
% ============================================================
\section{引言}

\subsection{物理动机}

理解核子内部的部分子结构是强子物理的核心问题之一。
胶子携带核子约 40--50\% 的动量，其部分子分布函数 (gluon PDF) $g(x,\mu)$
在高能散射过程中起着关键作用。
然而，胶子 PDF 的非微扰计算长期以来是一个巨大的挑战。

光锥 PDF 的原始定义为 Minkowski 空间中的光锥关联函数\cite{gluon_pdf_derivation}：
\begin{equation}\label{eq:lightcone_pdf}
x g(x, \mu) = \int \frac{dz^{-}}{2\pi P^{+}} e^{-i x P^{+} z^{-}}
    \langle P | F^{+\mu}_{a}(z^{-}) \, U_{ab}(z^{-}, 0) \,
    F^{b\mu}_{+}(0) (\mu) | P \rangle,
\end{equation}
其中 $U_{ab}(z^{-}, 0) = \mathcal{P}\exp[ig\int_{0}^{z^{-}} d\eta^{-} A^{+}(\eta^{-})]_{ab}$
为伴随表示中的 Wilson 线。
由于该定义涉及光锥时间 $z^{-}$，无法直接在 Euclidean 格点上计算。

\subsection{大动量有效理论 (LaMET)}

LaMET 由季向东于 2013 年提出\cite{lamet_original}，
通过在有限大动量 $P^{z}$ 下计算 Euclidean 空间的 quasi-PDF，
然后通过微扰匹配将其转换为光锥 PDF。
LaMET 因子化公式为\cite{gluon_pdf_derivation}：
\begin{equation}\label{eq:lamet_factorization}
x \tilde{g}(x, \Pz) = \int_{-\infty}^{\infty} \frac{dy}{|y|}
    C^{\text{hybrid}}\left(\frac{x}{y}, \lambda_s, \frac{\mu}{y\Pz}\right)
    \, y g(y, \mu) + \mathcal{O}\left(\frac{\QCD^{2}}{(x\Pz)^{2}},
    \frac{\QCD^{2}}{((1-x)\Pz)^{2}}\right).
\end{equation}

Quasi-PDF 由重整化后的格点矩阵元 Fourier 变换得到：
\begin{equation}\label{eq:quasi_pdf}
x \tilde{g}(x, \Pz) = \frac{1}{\mathcal{N}} \int \frac{dz}{2\pi\Pz}
    e^{-i x z \Pz} \, \hR(z, \Pz).
\end{equation}

\subsection{胶子PDF格点计算的关键挑战}

与夸克 PDF 不同，胶子 PDF 的格点计算面临以下独特困难：
\begin{enumerate}
    \item \textbf{非连通图 (Disconnected Diagrams)：}
        胶子算符不通过任何夸克传播子与核子连接，
        三点函数分解为核子两点函数与胶子 OPE 算符的乘积。

    \item \textbf{场强张量的格点构造：}
        胶子算符由场强张量 $F_{\mu\nu}$ 构建，
        需要在格点上通过 plaquette/Clover 叶构造。

    \item \textbf{Wilson 线的重整化：}
        非定域胶子算符中的空间 Wilson 线存在线性发散，
        需要特殊的混合方案重整化。

    \item \textbf{信噪比问题：}
        胶子算符的涨落远大于夸克算符，
        对统计精度和涂抹技术提出了更高要求。
\end{enumerate}

\subsection{本文结构}

第\ref{sec:physics}节建立完整的物理框架；
第\ref{sec:fmunu}节详述场强张量的格点构造及其代码实现；
第\ref{sec:ope}节阐述非定域胶子 OPE 算符的构造；
第\ref{sec:two_point}节介绍蒸馏法核子两点关联函数；
第\ref{sec:three_point}节推导非连通三点函数、比值方法与 Jackknife 重采样；
第\ref{sec:huangcl}节对 huangcl 的 \texttt{code.py} 进行完整的逐段分析；
第\ref{sec:workflow}节总结完整计算流程；
最后为附录，列出所有参考源。


% ============================================================
% 第二章：物理基础
% ============================================================
\section{物理基础：非极化胶子PDF的理论框架}\label{sec:physics}

\subsection{光锥胶子PDF}

光锥胶子PDF由 twist-2 光锥算符的 forward 矩阵元定义\cite{note_of_gluon_PDFs}。
在光锥规范 $A^{+}=0$ 下，非极化胶子 PDF 的最简形式为：
\begin{equation}\label{eq:gluon_pdf_lightcone}
x g(x) = \frac{1}{2} \int \frac{dz^{-}}{2\pi P^{+}} e^{-i x P^{+} z^{-}}
    \langle P | F^{+\mu}(z^{-}) F_{+\mu}(0) | P \rangle.
\end{equation}

\subsection{光锥关联函数的张量分解}

协变形式下，胶子矩阵元可参数化为\cite{note_of_gluon_PDFs, gluon_pdf_derivation}：
\begin{equation}\label{eq:M_munu}
M_{\mu\alpha; \lambda\beta}(z, p) = \langle p |
    F_{\mu\alpha}(z) \, W[z, 0] \, F_{\lambda\beta}(0) \, W[0, z] | p \rangle.
\end{equation}

通过 Lorentz 协变性，$M_{\mu\alpha;\lambda\beta}$ 可分解为六个独立的标量振幅函数。
对于非极化情况，twist-2 贡献来自不变振幅 $M_{pp}$：
\begin{equation}\label{eq:Mpp_relation}
g^{\alpha\beta} M_{+\alpha; \beta+}(z, p)
    = g^{ij} M_{+i; j+}(z, p)
    = -2 p_{+}^{2} \, M_{pp}(\nu, 0),
\end{equation}
其中 Ioffe 时间 $\nu = p_{3} z_{3}$。

\subsection{不变振幅与胶子PDF的关系}

$M_{pp}$ 是胶子 PDF 的 Fourier 变换\cite{note_of_gluon_PDFs}：
\begin{equation}\label{eq:Mpp_pdf}
- M_{pp}(\nu) = \frac{1}{2} \int_{-1}^{1} dx \, e^{-i x \nu} \, x g(x).
\end{equation}

这就是格点可计算量与连续光锥 PDF 之间的核心桥梁。

\subsection{非极化胶子算符的组合}

为了提取 $M_{pp}$，需要特定的 $M_{\mu\alpha;\lambda\beta}$ 线性组合。
在 Minkowski 空间中，非极化胶子流算符为\cite{note_of_gluon_PDFs}：
\begin{equation}\label{eq:Og_minkowski}
O^{(0)}_{g}(z) = \Big[ F^{0i}(z) W[z,0] F_{0i}(0) W[0,z]
    + 2 F^{12}(z) W[z,0] F_{12}(0) W[0,z] \Big]_{\text{Mink}}.
\end{equation}

利用 Minkowski--Euclidean 对应关系
$F^{(M)}_{0i} F^{(M)}_{0i} = -F^{(E)}_{0i} F^{(E)}_{0i}$ 和
$F^{(M)}_{ij} F^{(M)}_{ij} = F^{(E)}_{ij} F^{(E)}_{ij}$，
得到 Euclidean 空间的表达式\cite{gluon_pdf_derivation, 场强张量}：
\begin{equation}\label{eq:Og_euclidean}
\boxed{
O^{(0)}_{g}(z) = \Big[ -F_{0i}(z) W[z,0] F_{0i}(0) W[0,z]
    + 2 F_{12}(z) W[z,0] F_{12}(0) W[0,z] \Big]_{\text{Eucl}}.
}
\end{equation}

\subsection{非连通三点关联函数}

由于胶子算符不通过夸克线连接核子，
三点函数分解为非连通的乘积形式\cite{gluon_pdf_derivation}：
\begin{equation}\label{eq:C3pt_disc}
\boxed{
C_{\text{3pt}}(z, \Pz, \tsep, t)
    = \frac{1}{N_t} \sum_{t_0}
    \big\langle (O_g(z, t_0+t) - \langle O_g(z) \rangle)
    (C^{N}_{\text{2pt}}(\Pz, t_0+\tsep, t_0) - \langle C^{N}_{\text{2pt}} \rangle) \big\rangle.
}
\end{equation}

其中 $C^{N}_{\text{2pt}}$ 为核子两点关联函数，
$\langle \cdot \rangle$ 表示系综平均。
真空减除 (vacuum subtraction) 是消除非连通真空贡献的关键步骤。

\subsection{比值方法}

为消除源-汇重叠因子，构造三点与两点的比值\cite{gluon_pdf_derivation, 关联函数}：
\begin{equation}\label{eq:ratio_def}
\boxed{
R(\Pz; \tsep, t) \equiv \frac{C_{\text{3pt}}(\Pz; \tsep, t)}
    {C_{\text{2pt}}(\Pz; \tsep)}.
}
\end{equation}

在大时间分离极限下，激发态污染可通过两态截断描述：
\begin{equation}\label{eq:ratio_fit}
\boxed{
R(\Pz; \tsep, t) \approx c_0 + c_1 e^{-\Delta E (\tsep - t)}
    + c_1 e^{-\Delta E t},
}
\end{equation}
其中 $c_0$ 即为所需的裸 quasi-矩阵元 $\langle 0 | O_g(z) | 0 \rangle$，
$\Delta E = E_1 - E_0$ 为基态与第一激发态的能隙。


% ============================================================
% 第三章：场强张量
% ============================================================
\section{格点上的场强张量 $F_{\mu\nu}$}\label{sec:fmunu}

\subsection{连续统场强张量}

在连续统 Yang-Mills 理论中，场强张量定义为协变导数的对易子\cite{场强张量}：
\begin{equation}\label{eq:F_continuum}
F_{\mu\nu}(x) = -\frac{i}{g} [D_{\mu}(x), D_{\nu}(x)],
\qquad D_{\mu}(x) = \partial_{\mu} + i g A_{\mu}(x),
\end{equation}
其分量为：
\begin{equation}
F^{a}_{\mu\nu} = \partial_{\mu} A^{a}_{\nu}
    - \partial_{\nu} A^{a}_{\mu}
    - g f^{abc} A^{b}_{\mu} A^{c}_{\nu}.
\end{equation}

\subsection{链变量与Plaquette}

格点规范理论的基本自由度是链变量 (link variable)\cite{场强张量}：
\begin{equation}\label{eq:link}
U_{\mu}(x) = \exp(i a A_{\mu}(x)).
\end{equation}

最小的规范不变闭合回路是 plaquette (小 Wilson 环)：
\begin{equation}\label{eq:plaquette}
\boxed{
P_{\mu\nu}(x) = U_{\mu}(x) U_{\nu}(x+\hat{\mu})
    U^{\dagger}_{\mu}(x+\hat{\nu}) U^{\dagger}_{\nu}(x).
}
\end{equation}

通过 Baker-Campbell-Hausdorff (BCH) 展开到 $\mathcal{O}(a^{2})$：
\begin{equation}\label{eq:plaquette_bch}
P_{\mu\nu} = \exp\Big[ i a^{2} \big( \partial_{\mu} A_{\nu}
    - \partial_{\nu} A_{\mu} + i [A_{\mu}, A_{\nu}] \big)
    + \mathcal{O}(a^{3}) \Big]
    = 1 + i a^{2} F_{\mu\nu} + \mathcal{O}(a^{3}).
\end{equation}

\subsection{代码实现：Plaquette}

\begin{codeblock}{Operator.py --- plaquette\_new 函数}
def plaquette_new(gauge, mu, nu, Nt, Nx):
    """
    Compute P_{mu,nu}(x) = U_mu(x) U_nu(x+mu) U_mu^dag(x+nu) U_nu^dag(x)

    gauge shape: [t, z, y, x, dir, color, color]  (donghx convention)
    """
    pla = np.einsum(
        "tzyxab,tzyxbc->tzyxac",
        gauge[:, :, :, :, mu, :, :],
        np.roll(gauge, -1, 3 - mu)[:, :, :, :, nu, :, :],
    )
    #  U_mu(x) * U_nu(x+mu)

    pla = np.einsum(
        "tzyxab,tzyxcb->tzyxac",
        pla,
        np.roll(gauge, -1, 3 - nu)[:, :, :, :, mu, :, :].conj(),
    )
    #  * U_mu^dag(x+nu)

    pla = np.einsum(
        "tzyxab,tzyxcb->tzyxac",
        pla,
        gauge[:, :, :, :, nu, :, :].conj(),
    )
    #  * U_nu^dag(x)

    return pla
\end{codeblock}

\textbf{物理--代码对应：}
\begin{itemize}
    \item \texttt{gauge[:,:,:,:,mu,:,:]} $\longleftrightarrow$ $U_{\mu}(x)$
    \item \texttt{np.roll(gauge, -1, 3-mu)[:,:,:,:,nu,:,:]} $\longleftrightarrow$
        $U_{\nu}(x+\hat{\mu})$
    \item \texttt{np.roll(gauge, -1, 3-nu)[:,:,:,:,mu,:,:].conj()} $\longleftrightarrow$
        $U^{\dagger}_{\mu}(x+\hat{\nu})$
    \item \texttt{gauge[:,:,:,:,nu,:,:].conj()} $\longleftrightarrow$
        $U^{\dagger}_{\nu}(x)$
    \item \texttt{np.einsum("tzyxab,tzyxbc->tzyxac", ...)} $\longleftrightarrow$
        颜色空间的 $3\times 3$ 矩阵乘法
\end{itemize}

\subsection{Clover叶构造}

为了减小离散化误差并降低统计涨落，
采用 Clover 叶 (四个共面 plaquette 的平均)\cite{场强张量, gluon_pdf_derivation}：
\begin{equation}\label{eq:clover}
\boxed{
Q_{\mu\nu}(x) = P_{\mu\nu}(x)
    - P_{\nu,-\mu}(x)
    + P_{-\nu,\mu}(x)
    + P_{-\mu,-\nu}(x).
}
\end{equation}

四个 plaquette 分别对应：
\begin{itemize}
    \item $P_{\mu\nu}$ (右上)：$(\mu, \nu)$ 正方向 plaquette
    \item $-P_{\nu,-\mu}$ (左上)：$(\nu, -\mu)$ 方向 plaquette
    \item $+P_{-\nu,\mu}$ (左下)：$(-\nu, \mu)$ 方向 plaquette
    \item $+P_{-\mu,-\nu}$ (右下)：$(-\mu, -\nu)$ 方向 plaquette
\end{itemize}

Clover 叶的性质\cite{场强张量}：
\begin{itemize}
    \item 反对称性：$Q_{\nu\mu} = -Q_{\mu\nu}$
    \item 消除奇次离散化误差，领头阶误差为 $\mathcal{O}(a^{4})$
\end{itemize}

\subsection{Clover叶的代码实现}

\begin{codeblock}{Operator.py --- plaquette\_clover\_new 函数}
def plaquette_clover_new(gauge, mu, nu):
    """
    Compute clover term Q_{mu,nu} from four coplanar plaquettes.

    Returns anti-hermitian field strength operator:
    F_{mu,nu}^{clover} = -i/8 * (Q_{mu,nu} - Q_{nu,mu})
    """
    # 四个偏移的 plaquette
    gauge_rightup = gauge
    gauge_leftup  = np.roll(gauge, 1, 3 - mu)
    gauge_rightdown = np.roll(gauge, 1, 3 - nu)
    gauge_leftdown  = np.roll(gauge_leftup, 1, 3 - nu)

    # P_{mu,nu}  -- 右上
    pla_rightup = (compute 4-link products as in plaquette_new)

    # -P_{nu,-mu}  -- 左上 (负号来自方向反转)
    pla_leftup = (compute plaquette for (nu, -mu) direction)

    # +P_{-nu,mu}  -- 左下
    pla_leftdown = (compute plaquette for (-nu, mu) direction)

    # +P_{-mu,-nu}  -- 右下
    pla_rightdown = (compute plaquette for (-mu, -nu) direction)

    # 反对称化: Q - Q^dag
    ans = (pla_rightup  - pla_rightup.conj().transpose(0,1,2,3,5,4)
         + pla_leftup   - pla_leftup.conj().transpose(0,1,2,3,5,4)
         + pla_leftdown - pla_leftdown.conj().transpose(0,1,2,3,5,4)
         + pla_rightdown - pla_rightdown.conj().transpose(0,1,2,3,5,4))

    return -1j * ans / 8.0  # <-- F_{mu,nu} = -i/(8a^2) * (Q - Q^dag)
\end{codeblock}

\textbf{物理--代码对应：}
\begin{itemize}
    \item \texttt{pla\_rightup} $\longleftrightarrow$ $P_{\mu\nu}$
    \item \texttt{pla\_leftup} $\longleftrightarrow$ $-P_{\nu,-\mu}$
    \item \texttt{pla\_leftdown} $\longleftrightarrow$ $+P_{-\nu,\mu}$
    \item \texttt{pla\_rightdown} $\longleftrightarrow$ $+P_{-\mu,-\nu}$
    \item \texttt{X - X.conj().transpose(...5,4)} $\longleftrightarrow$
        $Q_{\mu\nu} - Q_{\nu\mu}$ (颜色矩阵反对称化)
    \item \texttt{-1j * ans / 8.0} $\longleftrightarrow$
        $-\frac{i}{8a^{2}}(Q_{\mu\nu} - Q_{\nu\mu})$ (因子 $1/a^{2}$ 在物理单位中恢复)
\end{itemize}

\subsection{全部 ($\mu$, $\nu$) 平面的Clover叶}

\begin{codeblock}{Operator.py --- plaquette\_clover\_all\_new 函数}
def plaquette_clover_all_new(gauge, Nt, Nx):
    """
    Compute F_{mu,nu} for all 4x4 (mu,nu) combinations.
    Returns shape: [4, 4, Nt, Nz, Ny, Nx, 3, 3]

    Diagonal (mu == nu) entries are zero-padded.
    """
    pla = np.zeros((4, 4, Nt, Nx, Nx, Nx, 3, 3), dtype=complex)
    for _mu in range(4):
        for _nu in range(4):
            if _mu == _nu:
                pla[_mu, _nu] = Temp_zeros   # F_{mu,mu} = 0
            else:
                pla[_mu, _nu] = plaquette_clover_new(gauge, _mu, _nu)
    return pla
\end{codeblock}

输出形状为 \texttt{[4, 4, Nt, Nz, Ny, Nx, 3, 3]}，前两个维度对应 $(\mu,\nu)$
索引。$\mu, \nu \in \{0,1,2,3\}$ 对应方向 $\{x, y, z, t\}$。

\subsection{对偶场强张量 $\tilde{F}_{\mu\nu}$}

对于部分胶子自旋算符的构造，需要用到对偶场强张量\cite{gluon_pdf_derivation}：
\begin{equation}
\tilde{F}^{\mu\nu} \equiv \frac{1}{2} \epsilon^{\mu\nu\rho\sigma} F_{\rho\sigma},
\end{equation}
其中 $\epsilon^{0123} = 1$。

代码中使用预先计算的 Levi-Civita 张量 \texttt{Tensor4}：
\begin{codeblock}{Operator.py --- plaquette\_clover\_all\_tilde 函数}
# Tensor4[i,j,k,l] = (1/2) * epsilon_{i,j,k,l}
# (zero if any indices coincide, +/-1/2 for even/odd permutations)
Tensor4 = np.zeros((4, 4, 4, 4), dtype=complex)
for i,j,k,l in itertools.product(range(4), repeat=4):
    if len({i,j,k,l}) != 4:  # 有重复指标 -> 0
        Tensor4[i,j,k,l] = 0
    else:
        sign = (+1 if parity of permutation is even else -1)
        Tensor4[i,j,k,l] = 0.5 * sign

def plaquette_clover_all_tilde(pla_all, Nt, Nx):
    pla_tilde = contract(
        "opmn,mntzyxab->optzyxab",  # 1/2 * epsilon contraction
        Tensor4,
        pla_all,
    )
    return pla_tilde
\end{codeblock}

\textbf{物理--代码对应：}
\begin{itemize}
    \item \texttt{Tensor4[o,p,m,n]} $\longleftrightarrow$
        $\frac{1}{2} \epsilon^{opmn}$
    \item \texttt{contract("opmn,mntzyxab->optzyxab", ...)} $\longleftrightarrow$
        $\tilde{F}^{op} = \frac{1}{2} \epsilon^{opmn} F_{mn}$
\end{itemize}


% ============================================================
% 第四章：非定域胶子OPE算符
% ============================================================
\section{非定域胶子OPE算符的构造}\label{sec:ope}

\subsection{格点上的OPE算符}

在格点上，非定域胶子 OPE 算符定义为\cite{gluon_pdf_derivation}：
\begin{equation}\label{eq:ope_lattice}
\boxed{
\mathcal{M}^{\mu\lambda, \nu\rho}(z)
    = \sum_{\vec{x}} \Tr \Big[ F^{\mu\lambda}(\vec{x} + z\hat{z}) \,
        U(\vec{x} + z\hat{z}, \vec{x}) \,
        F^{\nu\rho}(\vec{x}) \,
        U(\vec{x}, \vec{x} + z\hat{z}) \Big],
}
\end{equation}
其中 $U(\vec{x} + z\hat{z}, \vec{x})$ 是沿 $z$ 方向连接两个场强插入点的
空间 Wilson 线（由 $\hat{z}$ 方向的 link 变量乘积构成）。

\subsection{自旋平均的非极化算符}

非极化胶子 PDF 所需的乘性可重整化组合为\cite{gluon_pdf_derivation}：
\begin{equation}\label{eq:O_renormalizable}
\boxed{
\mathcal{O}(z) = \mathcal{M}^{tx; tx}(z)
    + \mathcal{M}^{ty; ty}(z)
    - 2 \mathcal{M}^{xy; xy}(z).
}
\end{equation}

这个特定组合保证了所有三项共享相同的单圈 UV 反常维度，
从而整个算符是乘性可重整化的（至少在单圈水平上）\cite{gluon_pdf_derivation}。

根据 Minkowski-Euclidean 对应关系：
\begin{align}
F^{(M)}_{ti} F^{(M)}_{ti} &= -F^{(E)}_{ti} F^{(E)}_{ti}, \\
F^{(M)}_{ij} F^{(M)}_{ij} &= +F^{(E)}_{ij} F^{(E)}_{ij}.
\end{align}

因此在 Euclidean 格点上\cite{note_of_gluon_PDFs, 场强张量}：
\begin{equation}\label{eq:O_euclidean_expanded}
\boxed{
\mathcal{O}(z)_{\text{Eucl}} \equiv
    -\mathcal{M}^{ti; ti}(z)
    + 2 \mathcal{M}^{ij; ij}(z).
}
\end{equation}

对 $z$ 方向为 $z$ 的情况，展开为：
\begin{equation}\label{eq:O_euclidean_components}
\mathcal{O}(z)_{\text{Eucl}} = -\mathcal{M}^{30;30}(z)
    - \mathcal{M}^{31;31}(z)
    + 2 \mathcal{M}^{01;01}(z).
\end{equation}

\subsection{OPE算符的MPI并行构造}

在 \texttt{Calc\_ope\_unpol.py} 中，MPI 的 4 个 rank 分别计算不同指标组合：

\begin{table}[H]
\centering
\caption{MPI rank 与 OPE 算符分量的对应关系 (Wilson线沿 $z$ 方向)}
\begin{tabular}{c|c|c|c|l}
\toprule
\textbf{Rank} & $\bm{\mu}$ & $\bm{\nu}$ & $\bm{(\mu,\nu)}$ 平面 & \textbf{物理意义} \\
\midrule
0 & 3 ($t$)  & 0 ($x$) & $t$-$x$ & $\mathcal{M}^{30;30}$ (电场 $E_x$) \\
1 & 3 ($t$)  & 1 ($y$) & $t$-$y$ & $\mathcal{M}^{31;31}$ (电场 $E_y$) \\
2 & 0 ($x$)  & 1 ($y$) & $x$-$y$ & $\mathcal{M}^{01;01}$ (磁场 $B_z$) \\
3 & --- & --- & --- & 空闲 \\
\bottomrule
\end{tabular}
\end{table}

\subsection{核心函数：operators\_new\_z0\_mu2}

\begin{codeblock}{Operator.py --- operators\_new\_z0\_mu2 函数 (简化)}
def operators_new_z0_mu2(gauge, zdir, pla, pla_tilde,
                          delta_z, mu, nu, mu2, nu2, Nt):
    """
    Compute O(z) = Tr[ F_{mu,nu}(z) * W^dag * Ftilde_{mu2,nu2}(0) * W ]

    zdir: Wilson line direction (0=x, 1=y, 2=z)
    pla:  F_{mu,nu} from clover construction
    pla_tilde: \tilde{F}_{mu,nu} = (1/2) * epsilon * F
    """
    z_dir = zdir

    # Step 1: Shift F_{mu,nu} from z=0 to z=delta_z
    ope = np.roll(pla[mu, nu], -delta_z, 3 - z_dir)
    # ope shape: [Nt, Nz, Ny, Nx, 3, 3]

    # Step 2: Backward Wilson line from z=delta_z to z=0
    #    W^dag(z,0) = U_z^dag(delta_z-1) * ... * U_z^dag(0)
    for _dz in range(delta_z):
        ope = np.einsum(
            "tzyxab,tzyxcb->tzyxac",
            ope,
            np.roll(gauge, -(delta_z - 1 - _dz), 3 - z_dir)
              [:, :, :, :, z_dir, :, :].conj(),
        )
    # 每个 einstein sum: ope * U_z^dag(step)  (颜色矩阵乘法)
    # 注意: .conj() = 厄米共轭 = U^dag

    # Step 3: Insert Ftilde_{mu2,nu2} at z=0
    ope = np.einsum(
        "tzyxab,tzyxbc->tzyxac",
        ope,
        pla_tilde[mu2, nu2],
    )
    # ope = ope * Ftilde(0)

    # Step 4: Forward Wilson line from z=0 back to z=delta_z
    #    W(0,z) = U_z(0) * U_z(1) * ... * U_z(delta_z-1)
    for _dz in range(delta_z):
        ope = np.einsum(
            "tzyxab,tzyxbc->tzyxac",
            ope,
            np.roll(gauge, -_dz, 3 - z_dir)
              [:, :, :, :, z_dir, :, :],
        )
    # 每个 einstein sum: ope * U_z(step)

    # Step 5: Color trace and spatial sum
    ans = np.trace(ope, axis1=4, axis2=5)  # Tr_color
    # ans shape: [Nt, Nz, Ny, Nx]
    op = np.sum(ans, axis=(1, 2, 3))       # sum over spatial volume
    # op shape: [Nt]

    return op
\end{codeblock}

\textbf{物理--代码对应的图示：}

代码中 operator 构造的图形表示为：
\begin{center}
\begin{tabular}{c}
$F_{\mu\nu}(z)$ at $z = \delta_z$ \\
$\downarrow$ \texttt{np.roll(pla[mu,nu], -delta\_z, ...)} \\
\framebox{$F_{\mu\nu}(z)$ shifted to $z=\delta_z$} \\
$\downarrow$ \texttt{for loop: ope * U\_z.conj()} \\
\framebox{向后Wilson线 $W^{\dagger}(z,0)$：从 $z=\delta_z$ 到 $0$} \\
$\downarrow$ \texttt{einsum: ope * pla\_tilde[mu2,nu2]} \\
\framebox{插入 $\tilde{F}_{\mu_2\nu_2}(0)$ at $z=0$} \\
$\downarrow$ \texttt{for loop: ope * U\_z} \\
\framebox{向前Wilson线 $W(0,z)$：从 $0$ 回到 $z=\delta_z$} \\
$\downarrow$ \texttt{np.trace(axis1=4, axis2=5)} \\
\framebox{颜色求迹 $\Tr_c[\cdots]$} \\
$\downarrow$ \texttt{np.sum(axis=(1,2,3))} \\
\framebox{空间求和 $\sum_{\vec{x}}$ → $\mathcal{O}(t; z)$}
\end{tabular}
\end{center}

\subsection{数据格式}

OPE 算符的输出格式为：
\begin{itemize}
    \item \textbf{文件名：} \texttt{ops\_mu\{mu\}\_nu\{nu\}\_dz\{delta\_z\}\_conf\{conf\_id\}.npy}
    \item \textbf{形状：} \texttt{[Nt, delta\_z]}（huancl加载后仅保留 \texttt{ops} 键的 \texttt{[Nt]} 数组）
    \item \textbf{含义：} 每个时间片 $t \in [0, N_t-1]$ 上，空间体积求和后的 OPE 算符值
\end{itemize}


% ============================================================
% 第五章：核子两点关联函数
% ============================================================
\section{核子两点关联函数：蒸馏法}\label{sec:two_point}

\subsection{蒸馏法概述}

蒸馏法 (distillation) 是一种将费米子场投影到三维规范协变 Laplacian 算符
的低本征模子空间上的涂抹方法\cite{distillation, 蒸馏方法}。
其核心优势在于：

\begin{enumerate}
    \item \textbf{低秩近似：} 使用 $N_{\text{ev}} \ll N_c \times V$ 个本征矢量，
        使所有传播子矩阵元可计算
    \item \textbf{保持对称性：} 涂抹算符与格点对称群对易
    \item \textbf{因子化结构：} 关联函数分解为 perambulator（传播子）和
        elemental（插值算符）的收缩
    \item \textbf{动量涂抹：} 可在蒸馏框架内自然实现 boosted 强子的动量涂抹
\end{enumerate}

\subsection{蒸馏算符}

三维规范协变 Laplacian 定义在时间片 $t$ 上\cite{蒸馏方法}：
\begin{equation}\label{eq:laplacian}
-\nabla^{2}_{xy}(t) = 6\delta_{xy}
    - \sum_{j=1}^{3} \big[ \tilde{U}_j(x,t) \delta_{x+\hat{j},y}
    + \tilde{U}^{\dagger}_j(x-\hat{j},t) \delta_{x-\hat{j},y} \big].
\end{equation}

求解本征值问题：
\begin{equation}
-\nabla^{2}(t) \, v^{(k)}(t) = \lambda_k(t) \, v^{(k)}(t),
\qquad v^{(k)\dagger} v^{(l)} = \delta_{kl}.
\end{equation}

蒸馏算符定义为最低 $N_{\text{ev}}$ 个本征矢量的投影\cite{蒸馏方法}：
\begin{equation}\label{eq:distillation_op}
\boxed{
\Box_{xy}(t) = \sum_{k=1}^{N_{\text{ev}}} v^{(k)}_x(t) v^{(k)\dagger}_y(t)
    \equiv V(t) V^{\dagger}(t).
}
\end{equation}

\subsection{Perambulator}

Perambulator 编码了蒸馏空间中的夸克传播\cite{蒸馏方法}：
\begin{equation}\label{eq:perambulator}
\boxed{
\tau^{(p,\bar{p})}_{\alpha\beta}(t', t) =
    V^{(p)\dagger}(t') \, M^{-1}_{\alpha\beta}(t', t) \, V^{(\bar{p})}(t),
}
\end{equation}
其中 $M^{-1}$ 是格点 Dirac 算符的逆，$\alpha,\beta$ 是 Dirac 旋量指标。

Perambulator 的形状为 \texttt{[Nt, 4, 4, Nev, Nev]}，分别对应
汇时间 $t'$、汇自旋、源自旋、汇本征模、源本征模。

\subsection{核子插值算符}

标准质子插值算符为\cite{关联函数}：
\begin{equation}\label{eq:nucleon_interp}
N_{\alpha}(\vec{P}; t) = \sum_{\vec{x}} e^{-i\vec{x}\cdot\vec{P}}
    \epsilon^{abc} (u^{T}_a C\gamma_5\gamma_t d_b) u_{c,\alpha}(\vec{x}; t).
\end{equation}

宇称投影算符 $P^{+} = \frac{1}{2}(1 + \gamma_t)$ 提取正宇称态。

\subsection{蒸馏法核子两点函数的Wick收缩}

在蒸馏框架下，核子两点函数的收缩由 VVV 张量和 perambulator 完成。
对于质子的 $uud$ 结构，有两个贡献项\cite{蒸馏方法}：

\begin{align}
C^{(2)}_{B}(t', t) &= \Phi^{(i,j,k)}(t')
    \, \tau^{(i,\bar{i})}_{d}(t',t)
    \, \tau^{(j,\bar{j})}_{u}(t',t)
    \, \tau^{(k,\bar{k})}_{u}(t',t)
    \, \Phi^{(\bar{i},\bar{j},\bar{k})*}(t) \\
    &\quad - \Phi^{(i,j,k)}(t')
    \, \tau^{(i,\bar{i})}_{d}(t',t)
    \, \tau^{(j,\bar{k})}_{u}(t',t)
    \, \tau^{(k,\bar{j})}_{u}(t',t)
    \, \Phi^{(\bar{i},\bar{j},\bar{k})*}(t). \nonumber
\end{align}

第一项是直接收缩，第二项来自 $u$ 夸克交换（费米子符号）。

\subsection{动量涂抹}

动量涂抹通过给本征矢量附加动量相关的相位实现\cite{distillation}：
\begin{equation}\label{eq:mom_smear}
\tilde{\xi}^{(k)}_{a}(\vec{z}, t) = e^{i\vec{\zeta}\cdot\vec{z}}
    \, \xi^{(k)}_{a}(\vec{z}, t),
\end{equation}
其中涂抹动量 $\vec{\zeta} = \zeta \times (2\pi/L)\hat{z}$，
$\zeta$ 为整数涂抹相位参数。

\subsection{2pt数据格式}

huangcl 代码加载的 2pt 数据来自 donghx 的蒸馏计算：
\begin{itemize}
    \item \textbf{文件名：}
        \texttt{twopt\_slice\_pp\_Px0Py0Pz2\_eginphase2\_Cg5g4\_nopol\_ss\_conf\{conf\_id\}.npy}
    \item \textbf{形状：} \texttt{[Nt, Nt]} $=$ \texttt{[t\_snk, t\_src]}
    \item \textbf{含义：} 正宇称投影后的核子两点关联函数
        $C_{\text{2pt}}(t_{\text{snk}}, t_{\text{src}}; \vec{P})$
    \item \textbf{参数：} $P_z = 2$ (动量), moms smear phase = 2,
        插值器 \texttt{Cg5g4} $= C\gamma_5\gamma_4$
\end{itemize}


% ============================================================
% 第六章：非连通三点关联函数与比值方法
% ============================================================
\section{非连通三点关联函数与比值方法}\label{sec:three_point}

\subsection{非连通三点函数的构造}

回顾公式 (\ref{eq:C3pt_disc})：
\begin{equation*}
C_{\text{3pt}}(z, \Pz, \tsep, t) =
    \frac{1}{N_t} \sum_{t_0}
    \big\langle \delta O_g(z, t_0+t) \;
    \delta C^{N}_{\text{2pt}}(\Pz, t_0+\tsep, t_0) \big\rangle,
\end{equation*}
其中 $\delta X \equiv X - \langle X \rangle$ 表示真空减除。

\subsection{huangcl代码中的数据结构}

代码加载以下数据：
\begin{enumerate}
    \item \textbf{2pt 关联函数} \texttt{\_corr[Nconf, Nt, Nt]}：
        $C_{\text{2pt}}(t_{\text{snk}}, t_{\text{src}})$

    \item \textbf{OPE 算符分量}（各 \texttt{[Nconf, Nx, Nt]}）：
        \begin{itemize}
            \item \texttt{\_ope\_01}：$\mathcal{O}_{\mu=0,\nu=1}(t)$
                -- $F_{01}$ (磁场 $B_z$) 的双线性
            \item \texttt{\_ope\_30}：$\mathcal{O}_{\mu=3,\nu=0}(t)$
                -- $F_{30}$ (电场 $E_x$) 的双线性
            \item \texttt{\_ope\_31}：$\mathcal{O}_{\mu=3,\nu=1}(t)$
                -- $F_{31}$ (电场 $E_y$) 的双线性
        \end{itemize}
\end{enumerate}

\subsection{OPE线性组合}

代码第89行执行：
\begin{codeblock}{code.py --- OPE 组合 (line 89)}
_ope = -_ope_30 - _ope_31 + 2 * _ope_01
\end{codeblock}

这精确对应公式 (\ref{eq:O_euclidean_components})：
\begin{equation*}
\mathcal{O}(z)_{\text{Eucl}} = -\mathcal{M}^{30;30}(z)
    - \mathcal{M}^{31;31}(z)
    + 2 \mathcal{M}^{01;01}(z).
\end{equation*}

\textbf{符号验证：}
\begin{itemize}
    \item Minkowski 非极化算符：
        $O_g^{(0)} = F^{0i}F_{0i} + 2F^{12}F_{12}$ （两项均为正）
    \item Wick 转动：$F^{0i}_{(M)}F^{(M)}_{0i} = -F^{0i}_{(E)}F^{(E)}_{0i}$，
        故 $F^{0i}_{(M)}F^{(M)}_{0i} \to -F^{0i}_{(E)}F^{(E)}_{0i}$
    \item 因此 Euclidean：$O_g^{(0)} = -F_{30}F_{30} - F_{31}F_{31} + 2F_{01}F_{01}$
    \item 代码：\texttt{- \_ope\_30 - \_ope\_31 + 2 * \_ope\_01} ✓
\end{itemize}

\subsection{时间平移与相对时间表示}

关联函数需要从绝对时间表示转换为相对时间表示。
定义源时间 $t_i$，汇时间 $t_i + \Delta t$，算符插入时间 $t_i + \Delta\tau$：

\begin{codeblock}{code.py --- 时间平移 (lines 92--104)}
max_t = 20  # maximum source-sink separation

# 初始化相对时间数组
_corr2_rel = np.zeros((Nconf, Nt, max_t), dtype=complex)
#  [conf, ti(source time), dt = t_snk - t_src]

_ope_rel = np.zeros((Nconf, Nt, max_t, Nx), dtype=complex)
#  [conf, ti, dtau = t_ins - t_src, z]

for ti in range(Nt):
    # 2pt: 从绝对时间 (t_snk, t_src=ti) 转换为相对时间 (ti, dt)
    corr2_shift = np.roll(_corr[:, :, ti], shift=-ti, axis=1)
    _corr2_rel[:, ti, :] = corr2_shift[:, :max_t]

    # OPE: 从绝对时间 (t, z) 转换为相对时间 (ti, dtau, z)
    ope_shift = np.roll(_ope, shift=-ti, axis=1)
    _ope_rel[:, ti, :, :] = ope_shift[:, :max_t, :]
\end{codeblock}

\textbf{物理--代码对应：}
\begin{itemize}
    \item \texttt{ti} $\longleftrightarrow$ 源时间 $t_0$
    \item \texttt{dt} $\longleftrightarrow$ 源-汇分离 $t_{\text{sep}} = t_{\text{snk}} - t_0$
    \item \texttt{dtau} $\longleftrightarrow$ 算符插入时间 $t = t_{\text{ins}} - t_0$
    \item \texttt{np.roll(..., shift=-ti)} $\longleftrightarrow$
        时间平移，使 $t_0$ 对齐到 $t=0$
    \item \texttt{[:max\_t]} $\longleftrightarrow$ 最大时间分离截断
\end{itemize}

\subsection{非连通三点函数的构造}

\begin{codeblock}{code.py --- 3pt 构造 (lines 106--114)}
_corr3 = np.zeros((Nconf, Nt, max_t, max_t, Nx), dtype=complex)
# [conf, ti, dt, dtau, z]

for dt in range(max_t):       # source-sink separation
    for dtau in range(dt + 1): # operator insertion time
        c2_slice = _corr2_rel[:, :, dt]
        # shape: [conf, ti]  -- C_2pt at fixed dt for all ti

        ope_slice = _ope_rel[:, :, dtau, :]
        # shape: [conf, ti, z]  -- OPE at fixed dtau for all ti

        c3_instance = ope_slice * c2_slice[:, :, np.newaxis]
        # shape: [conf, ti, z]  -- O_g(z, ti+dtau) * C_2pt(ti+dt, ti)

        _corr3[:, :, dt, dtau, :] = c3_instance
\end{codeblock}

\textbf{物理--代码对应：}
\begin{itemize}
    \item \texttt{c2\_slice[:,:,dt]} $\longleftrightarrow$
        $C_{\text{2pt}}(t_0 + \Delta t, t_0; \vec{P})$
    \item \texttt{ope\_slice[:,:,dtau,:]} $\longleftrightarrow$
        $\mathcal{O}_g(z, t_0 + \Delta\tau)$
    \item \texttt{c3\_instance = ope\_slice * c2\_slice[:,:,np.newaxis]} $\longleftrightarrow$
        $\mathcal{O}_g(z, t_0+\Delta\tau) \times C_{\text{2pt}}(t_0+\Delta t, t_0)$
    \item 这是公式 (\ref{eq:C3pt_disc}) 的被积函数
        （在求平均和真空减除之前）
\end{itemize}

\subsection{重采样：Jackknife方法}

\begin{codeblock}{code.py --- resample 函数 (lines 53--64)}
def resample(corr, jack, Nsample):
    n_conf = corr.shape[0]
    if jack:
        # Jackknife: 依次省略每个组态
        re_corr = (n_conf * corr.mean(0) - corr) / (n_conf - 1)
    else:
        # Bootstrap: 有放回随机重采样
        rng = np.random.default_rng(seed=seed)
        idx = rng.integers(0, n_conf, size=(Nsample, n_conf))
        re_corr = corr[idx].mean(1)
    return re_corr
\end{codeblock}

\textbf{Jackknife 公式：}
\begin{equation}\label{eq:jackknife}
\boxed{
C^{(i)}_{\text{JK}} = \frac{N C_{\text{full}} - C_i}{N - 1},
}
\end{equation}
其中 $C_{\text{full}} = \frac{1}{N} \sum_i C_i$，
$C_i$ 为第 $i$ 个组态上的观测量，
$C^{(i)}_{\text{JK}}$ 为省略第 $i$ 个组态后的 Jackknife 样本。

Jackknife 误差公式：
\begin{equation}\label{eq:jackknife_error}
\boxed{
\sigma_{\text{JK}} = \sqrt{\frac{N-1}{N}
    \sum_{i=1}^{N} (C^{(i)}_{\text{JK}} - \bar{C}_{\text{JK}})^2},
}
\end{equation}
其中 $\bar{C}_{\text{JK}} = \frac{1}{N} \sum_i C^{(i)}_{\text{JK}}$。

代码中 \texttt{sem} 函数实现了这一误差计算：
\begin{codeblock}{code.py --- sem 函数 (lines 46--50)}
def sem(data, jack):
    error = data.std(0)
    if jack:
        error = error * np.sqrt(data.shape[0] - 1)
    return error
\end{codeblock}

\textbf{推导：} 对于 Jackknife 样本 $\{C^{(i)}_{\text{JK}}\}$，
其标准差与原始样本均值标准差的关系为：
\begin{equation*}
\sigma_{\text{JK}} = \sigma_{\text{sample}} \cdot \sqrt{N-1},
\end{equation*}
其中 $\sigma_{\text{sample}} = \texttt{data.std(0)}$。
因此 \texttt{error * np.sqrt(N-1)} 给出正确的 Jackknife 误差。

\subsection{非连通减除与比值}

\begin{codeblock}{code.py --- 非连通减除与比值 (lines 127--132)}
# 非连通减除：减去 OPE 平均值 × 2pt 平均值 (真空贡献)
corr3_disc = (
    corr3 - corr2[:, :, :, np.newaxis, np.newaxis]
          * ope[:, :, np.newaxis, :, :]
)
# shape: [Nsample, ti, dt, dtau, z]

# 比值：对 ti 求平均
ratio = np.mean(
    (corr3_disc / corr2[:, :, :, np.newaxis, np.newaxis]),
    axis=1
)
# shape: [Nsample, dt, dtau, z]
\end{codeblock}

\textbf{物理--代码对应：}

\begin{enumerate}
    \item \texttt{corr3\_disc} $\longleftrightarrow$ 公式 (\ref{eq:C3pt_disc}) 的分子：
        \begin{equation*}
        C^{\text{disc}}_{\text{3pt}} = \langle O_g \cdot C_{\text{2pt}} \rangle
            - \langle O_g \rangle \cdot \langle C_{\text{2pt}} \rangle
        \end{equation*}
        这正是连接部分：$\langle (O_g - \langle O_g \rangle)
        (C_{\text{2pt}} - \langle C_{\text{2pt}} \rangle) \rangle$

    \item \texttt{ratio} $\longleftrightarrow$ 公式 (\ref{eq:ratio_def})：
        \begin{equation*}
        R(\Delta t, \Delta\tau, z) = \frac{1}{N_t} \sum_{t_i}
            \frac{C^{\text{disc}}_{\text{3pt}}(t_i; \Delta t, \Delta\tau, z)}
                 {C_{\text{2pt}}(t_i; \Delta t)}
        \end{equation*}
        对源时间 $t_i$ 的平均利用了时间平移不变性。
\end{enumerate}


% ============================================================
% 第七章：huangcl代码完整分析
% ============================================================
\section{huangcl \texttt{code.py} 完整剖析}\label{sec:huangcl}

\subsection{代码结构总览}

huangcl 的 \texttt{code.py}（约 219 行）实现了一个完整的
非极化胶子 PDF 分析流程。其结构如下：

\begin{enumerate}
    \item \textbf{配置参数} (lines 1--24)：格点参数、动量、统计方法
    \item \textbf{数据分配} (lines 39--43)：初始化存储数组
    \item \textbf{统计工具函数} (lines 46--64)：sem() 和 resample()
    \item \textbf{数据加载} (lines 71--84)：加载 2pt 和 OPE 数据
    \item \textbf{OPE 组合} (line 89)：线性组合为总 OPE 算符
    \item \textbf{时间平移} (lines 92--104)：绝对时间→相对时间
    \item \textbf{3pt 构造} (lines 106--114)：2pt 与 OPE 的乘积
    \item \textbf{内存清理} (lines 116--117, 124--125, 134--135)
    \item \textbf{重采样} (lines 120--122)：Jackknife
    \item \textbf{非连通减除与比值} (lines 127--132)
    \item \textbf{绘图} (lines 147--207)：ratio vs $\Delta\tau - \Delta t/2$
    \item \textbf{性能监控} (lines 210--216)：峰值内存
\end{enumerate}

\subsection{格点参数}

\begin{codeblock}{code.py --- 参数配置 (lines 11--24)}
conf_name = "beta6.20_mu-0.2770_ms-0.2400_L24x72"
conf_short = "L24x72"
Nconf = 200    # 组态数量
Nt = 72        # 时间方向格点数
Nx = 24        # 空间方向格点数

Px = 0         # x方向动量 (2π/L 单位)
Py = 0         # y方向动量
Pz = 2         # z方向动量 (P^z = 2 × 2π/L)

test = False   # 测试模式 (仅3个组态)
jack = True    # 使用 Jackknife 重采样
Nsample = 3000 # Bootstrap样本数 (仅 jack=False 时使用)
\end{codeblock}

\textbf{物理参数：}
\begin{itemize}
    \item \textbf{规范系综：} $\beta = 6.20$, $m_u = m_d = -0.2770$, $m_s = -0.2400$,
        $24^3 \times 72$ 格点
    \item \textbf{格距：} $a \approx 0.105$ fm
    \item \textbf{核子动量：} $P_z = 2 \times 2\pi/(24a) \approx 0.49$ GeV
    \item \textbf{组态数：} 200 个组态，间隔 200 轨迹（组态 ID: 6250, 6450, ..., 46050）
\end{itemize}

\subsection{数据加载详解}

\begin{codeblock}{code.py --- 数据加载 (lines 71--84)}
for i in range(Nconf):
    conf_id = 6250 + i * 200  # 组态ID: 6250--46050, 步长200

    # 2pt 关联函数: C_2pt(t_snk, t_src)
    _corr[i] = np.load(
        f".../2pt_Result/{conf_name}/momsmear2z/{conf_id}/"
        f"twopt_slice_pp_Px{Px}Py{Py}Pz{Pz}_eginphase2"
        f"_Cg5g4_nopol_ss_conf{conf_id}.npy"
    )
    # shape: [Nt, Nt] -- [t_snk, t_src]

    # OPE: mu=0, nu=1 (F_01, 磁场 B_z 分量)
    _ope_01[i] = np.load(
        f".../Ope_Gluon/Result_hpy_4D_10times/{conf_short}/zdir/{conf_id}/"
        f"ops_mu0_nu1_dz24_conf{conf_id}.npz"
    )["ops"]
    # shape: [Nx, Nt] -- [z, t]

    # OPE: mu=3, nu=0 (F_30, 电场 E_x 分量)
    _ope_30[i] = np.load(
        f".../Ope_Gluon/Result_hpy_4D_10times/{conf_short}/zdir/{conf_id}/"
        f"ops_mu3_nu0_dz24_conf{conf_id}.npz"
    )["ops"]

    # OPE: mu=3, nu=1 (F_31, 电场 E_y 分量)
    _ope_31[i] = np.load(
        f".../Ope_Gluon/Result_hpy_4D_10times/{conf_short}/zdir/{conf_id}/"
        f"ops_mu3_nu1_dz24_conf{conf_id}.npz"
    )["ops"]
\end{codeblock}

\textbf{数据来源对应关系：}

\begin{table}[H]
\centering
\caption{huangcl 加载的数据与其生成代码的对应关系}
\begin{tabular}{c|c|l}
\toprule
\textbf{数据} & \textbf{形状} & \textbf{生成代码} \\
\midrule
2pt 关联函数 & \texttt{[Nt, Nt]} &
    \texttt{donghx/2pt\_proton\_Cg5gmu\_L24x72\_mom2\_zdir\_gpu.py} \\
OPE $\mu=0,\nu=1$ & \texttt{[Nx, Nt]} &
    \texttt{donghx/Calc\_ope\_unpol.py} (rank 2)\\
OPE $\mu=3,\nu=0$ & \texttt{[Nx, Nt]} &
    \texttt{donghx/Calc\_ope\_unpol.py} (rank 0)\\
OPE $\mu=3,\nu=1$ & \texttt{[Nx, Nt]} &
    \texttt{donghx/Calc\_ope\_unpol.py} (rank 1)\\
\bottomrule
\end{tabular}
\end{table}

\subsection{完整的数据流图}

下图总结了从原始数据到最终 ratio 的完整张量流：

\begin{center}
\fbox{
\begin{minipage}{0.92\textwidth}
\textbf{输入数据：}
\begin{itemize}
    \item \texttt{\_corr} \texttt{[Nconf, Nt, Nt]}：$C_{\text{2pt}}(t_{\text{snk}}, t_{\text{src}})$
    \item \texttt{\_ope\_\{01,30,31\}} \texttt{[Nconf, Nx, Nt]}：
        $\mathcal{O}_{\mu\nu}(z, t)$
\end{itemize}

\bigskip
\hrule
\bigskip

\textbf{步骤 1 --- OPE 组合 (line 89)：}
\begin{itemize}
    \item \texttt{\_ope = -\_ope\_30 - \_ope\_31 + 2 * \_ope\_01}
    \item \texttt{\_ope} \texttt{[Nconf, Nt, Nx]} (转置后)：
        $\mathcal{O}_g(z, t) = -\mathcal{M}^{30;30} - \mathcal{M}^{31;31} + 2\mathcal{M}^{01;01}$
\end{itemize}

\bigskip
\hrule
\bigskip

\textbf{步骤 2 --- 时间平移 (lines 99--104)：}
\begin{itemize}
    \item \texttt{\_corr2\_rel[:, ti, dt]} =
        $C_{\text{2pt}}(t_i + \Delta t, t_i)$
    \item \texttt{\_ope\_rel[:, ti, dtau, z]} =
        $\mathcal{O}_g(z, t_i + \Delta\tau)$
\end{itemize}

\bigskip
\hrule
\bigskip

\textbf{步骤 3 --- 3pt 构造 (lines 109--114)：}
\begin{itemize}
    \item \texttt{\_corr3[:, ti, dt, dtau, z]} =
        $\mathcal{O}_g(z, t_i+\Delta\tau) \times C_{\text{2pt}}(t_i+\Delta t, t_i)$
\end{itemize}

\bigskip
\hrule
\bigskip

\textbf{步骤 4 --- Jackknife (lines 120--122)：}
\begin{itemize}
    \item 对 \texttt{\_corr2\_rel}, \texttt{\_ope\_rel}, \texttt{\_corr3} 分别重采样
\end{itemize}

\bigskip
\hrule
\bigskip

\textbf{步骤 5 --- 非连通减除与比值 (lines 127--132)：}
\begin{itemize}
    \item \texttt{corr3\_disc = corr3 - corr2 * ope}
    \item \texttt{ratio = mean(corr3\_disc / corr2, axis=1)} (对 $t_i$ 平均)
    \item \texttt{ratio} \texttt{[Nsample, dt, dtau, z]}：
        $R(\Delta t, \Delta\tau, z)$
\end{itemize}
\end{minipage}
}
\end{center}

\subsection{绘图分析}

\begin{codeblock}{code.py --- 绘图段 (lines 147--207)}
target_z = 2          # 选择 z = 2 (Wilson线长度为2a)
dt_list = [4,5,6,7,8,9,10]  # 源-汇分离

ratio_mean = ratio.mean(0)  # 对 Jackknife 样本求平均
ratio_err  = sem(ratio, jack)

for i, dt in enumerate(dt_list):
    tau_vals = np.arange(0, dt + 1)

    # 横坐标: tau - dt/2  (算子插入时间相对于中点)
    x_vals = tau_vals - dt / 2.0

    # 纵坐标: R(dt, dtau, z=2)
    y_vals = ratio_mean[dt, tau_vals, target_z]
    y_errs = ratio_err[dt, tau_vals, target_z]

    ax.errorbar(x_vals, y_vals, yerr=y_errs,
                fmt="x", label=f"z={target_z}, tsep={dt}")
\end{codeblock}

\textbf{物理意义：}
\begin{itemize}
    \item \textbf{横坐标：} $\Delta\tau - \Delta t/2$ —— 算符插入时间相对于源-汇中点的偏移。
        当 $\Delta\tau = \Delta t/2$（中点）时，激发态污染最小。
    \item \textbf{纵坐标：} $R(\Delta t, \Delta\tau, z)$ ——
        在特定 Wilson 线长度 $z$ 下的裸矩阵元估计。
    \item \textbf{不同曲线：} 对应不同的源-汇分离 $\Delta t$。
        在 $\Delta t \to \infty$ 极限下，所有曲线应趋于相同的平台值 $c_0$。
    \item \textbf{平台行为：} 在 $\Delta\tau - \Delta t/2 \approx 0$ 附近，
        $R$ 应该对 $\Delta\tau$ 呈现平台，表明基态饱和。
\end{itemize}

\subsection{内存管理}

代码中有三次显式的 \texttt{del} + \texttt{gc.collect()} 操作：

\begin{codeblock}{code.py --- 内存管理 (代表性段)}
# 第一次清理 (lines 116--117):
# 三维数组 _corr[Nconf,Nt,Nt] 和四个 OPE 分量不再需要
del _corr, _ope, _ope_30, _ope_31, _ope_01
gc.collect()

# 第二次清理 (lines 124--125):
# 原始数据已被重采样，清除未重采样的数组
del _corr2_rel, _ope_rel, _corr3
gc.collect()

# 第三次清理 (lines 134--135):
# 重采样后的 corr3, corr2, ope 不再需要
del corr3, corr2, ope, corr3_disc
gc.collect()
\end{codeblock}

\textbf{内存估算（每个组态阶段）：}
\begin{itemize}
    \item \texttt{\_corr}: $200 \times 72 \times 72 \times 16$ bytes $\approx 16.6$ MB
    \item \texttt{\_corr3}: $200 \times 72 \times 20 \times 20 \times 24 \times 16$ bytes
        $\approx 2.2$ GB
\end{itemize}

这解释了为什么代码需要在中间步骤清理内存——\texttt{\_corr3} 占据了大量内存。


% ============================================================
% 第八章：完整计算流程
% ============================================================
\section{完整计算工作流}\label{sec:workflow}

\subsection{端到端工作流}

非极化胶子 PDF 的完整格点计算流程如图\ref{fig:workflow}所示。

\begin{figure}[H]
\centering
\fbox{
\begin{minipage}{0.92\textwidth}
\small
\textbf{Phase 1: 数据生成 (donghx 代码, 在GPU集群上运行)}

\begin{enumerate}
    \item 读取 ILDG 格式的规范组态 $\{U_\mu(x)\}$
    \item 构造 Clover 叶 $\to F_{\mu\nu}(x)$ (所有 $(\mu,\nu)$ 平面)
    \item 计算对偶场强 $\tilde{F}_{\mu\nu} = \frac{1}{2}\epsilon_{\mu\nu\rho\sigma}F^{\rho\sigma}$
    \item 构造非定域 OPE 算符 $\mathcal{O}_{\mu\nu}(z,t)$
        (含 Wilson 线 $W[z,0]$)
    \item MPI 并行：rank 0/1/2 分别计算三组 $(\mu,\nu)$ 分量
    \item 输出：\texttt{ops\_mu\{mu\}\_nu\{nu\}\_dz24\_conf\{id\}.npz}
\end{enumerate}

\medskip
\textbf{Phase 2: 蒸馏计算 (donghx 代码, 在GPU集群上运行)}

\begin{enumerate}
    \item 求解 3D Laplacian 本征值问题 $\to$ $\{v^{(k)}(t)\}_{k=1}^{N_{\text{ev}}}$
    \item 动量涂抹本征矢量 $\tilde{v}^{(k)}(\vec{x},t) = e^{i\vec{\zeta}\cdot\vec{x}} v^{(k)}(\vec{x},t)$
    \item 计算 perambulator $\tau(t',t)$ (夸克传播子在蒸馏空间的投影)
    \item 计算 VVV 张量 (核子算符的蒸馏构造)
    \item Wick 收缩 + 宇称投影 $\to C_{\text{2pt}}(t_{\text{snk}}, t_{\text{src}})$
    \item 输出：\texttt{twopt\_slice\_pp\_...\_nopol\_ss\_conf\{id\}.npy}
\end{enumerate}

\medskip
\textbf{Phase 3: 分析 (huangcl \texttt{code.py}, 在CPU节点上运行)}

\begin{enumerate}
    \item 加载所有组态的 2pt 和 OPE 数据
    \item OPE 线性组合 $\mathcal{O}_g = -\mathcal{O}_{30} - \mathcal{O}_{31} + 2\mathcal{O}_{01}$
    \item 绝对时间 $\to$ 相对时间
    \item 构造非连通 3pt：
        $C_{\text{3pt}} = \mathcal{O}_g(z, t_i+\Delta\tau) \times C_{\text{2pt}}(t_i+\Delta t, t_i)$
    \item Jackknife 重采样
    \item 真空减除 + 比值：
        $R = \langle (C_{\text{3pt}} - \langle C_{\text{3pt}}\rangle) / C_{\text{2pt}} \rangle_{t_i}$
    \item 绘制 ratio 图：$R(z, \Delta t, \Delta\tau)$ vs $\Delta\tau - \Delta t/2$
\end{enumerate}

\medskip
\textbf{Phase 4: 后续分析 (未包含在 huangcl 代码中)}

\begin{enumerate}
    \item 两态拟合 $R \approx c_0 + c_1 e^{-\Delta E(\Delta t - \Delta\tau)} + c_1 e^{-\Delta E \Delta\tau}$
    \item 提取 $c_0(z)$ = 裸 quasi-矩阵元
    \item 混合方案重整化
    \item Fourier 变换 $\to$ quasi-PDF $\tilde{g}(x, \Pz)$
    \item LaMET 匹配 $\to$ 光锥 PDF $g(x, \mu)$
    \item 连续极限 + 无限动量外推
\end{enumerate}
\end{minipage}
}
\caption{非极化胶子 PDF 的完整格点计算工作流}
\label{fig:workflow}
\end{figure}

\subsection{关键物理量总结}

\begin{table}[H]
\centering
\caption{代码中的核心物理量与 Python 数组的对应关系}
\begin{tabular}{c|c|l}
\toprule
\textbf{符号} & \textbf{Python 数组} & \textbf{含义} \\
\midrule
$U_\mu(x)$ & \texttt{gauge[t,z,y,x,mu,c1,c2]} & 规范链变量 \\
$F_{\mu\nu}(x)$ & \texttt{pla[mu,nu,t,z,y,x,c1,c2]} & 场强张量 (Clover) \\
$\tilde{F}_{\mu\nu}(x)$ & \texttt{pla\_tilde[mu,nu,...]} & 对偶场强张量 \\
$\mathcal{O}_{\mu\nu}(z,t)$ & \texttt{\_ope\_\{mn\}[conf,z,t]} & OPE 算符分量 \\
$\mathcal{O}_g(z,t)$ & \texttt{\_ope[conf,t,z]} & 总 OPE 算符 \\
$C_{\text{2pt}}(t_s, t_0)$ & \texttt{\_corr[conf,t\_snk,t\_src]} & 核子两点函数 \\
$C_{\text{2pt}}(t_i; \Delta t)$ & \texttt{\_corr2\_rel[conf,ti,dt]} & 相对时间 2pt \\
$\mathcal{O}_g(t_i; z, \Delta\tau)$ & \texttt{\_ope\_rel[conf,ti,dtau,z]} & 相对时间 OPE \\
$C_{\text{3pt}}(t_i; \Delta t, \Delta\tau, z)$ & \texttt{\_corr3[conf,ti,dt,dtau,z]} & 非连通 3pt \\
$C^{\text{disc}}_{\text{3pt}}$ & \texttt{corr3\_disc[sample,ti,dt,dtau,z]} & 真空减除后 3pt \\
$R(\Delta t, \Delta\tau, z)$ & \texttt{ratio[sample,dt,dtau,z]} & 比值 \\
\bottomrule
\end{tabular}
\end{table}


% ============================================================
% 第九章：总结
% ============================================================
\section{总结}\label{sec:summary}

本文系统地分析了 huangcl 编写的非极化胶子 PDF 格点计算代码
(\texttt{examples/huangcl/code.py})，
建立了从连续统 QCD 物理公式到具体 Python 实现的完整对应关系。

\textbf{主要要点：}

\begin{enumerate}
    \item \textbf{场强张量的格点构造：}
        Plaquette $\to$ Clover 叶 $\to$ $F_{\mu\nu}$ 的离散化路径，
        其核心代码在 \texttt{Operator.py} 中实现。
        Clover 叶通过四个共面 plaquette 的反对称化组合，
        实现了 $\mathcal{O}(a^2)$ 的离散化精度。

    \item \textbf{非定域胶子算符：}
        $\mathcal{M}^{\mu\lambda,\nu\rho}(z) = \Tr[F^{\mu\lambda}(z) W F^{\nu\rho}(0) W^{\dagger}]$
        在代码中通过 \texttt{operators\_new\_z0\_mu2} 函数实现，
        依次执行：(1) 场强偏移、(2) 向后 Wilson 线、
        (3) 对偶场强插入、(4) 向前 Wilson 线、
        (5) 颜色求迹 + 空间求和。

    \item \textbf{非极化算符组合：}
        代码中 \texttt{- \_ope\_30 - \_ope\_31 + 2 * \_ope\_01}
        精确对应 Euclidean 空间的非极化胶子流
        $-F_{30}F_{30} - F_{31}F_{31} + 2F_{01}F_{01}$，
        其符号源于 Minkowski-Euclidean 对应中的 Wick 转动效应。

    \item \textbf{非连通三点函数：}
        代码通过 $\mathcal{O}_g \times C_{\text{2pt}}$ 的乘积
        构造非连通三点函数，随后进行 Jackknife 重采样、
        真空减除和比值计算，最终得到裸矩阵元的统计估计。

    \item \textbf{完整工作流：}
        donghx 的 GPU 代码生成 OPE 算符和蒸馏 2pt 数据
        (Phase 1--2)，
        huangcl 的 CPU 代码加载这些数据并完成
        比值分析和可视化 (Phase 3)，
        后续的两态拟合、重整化、Fourier 变换和 LaMET 匹配
        在 zhangxin 的工作流代码中完成 (Phase 4)。

    \item \textbf{统计方法：}
        Jackknife 重采样正确处理了组态间的关联，
        代码中 \texttt{sem} 函数的 $\sqrt{N-1}$ 因子
        将重采样样本的标准差转换为原始样本的 Jackknife 误差。
\end{enumerate}

huangcl 的代码虽然紧凑（约 220 行），但完整实现了
从原始格点数据到裸 quasi-矩阵元比值 $R(\Delta t, \Delta\tau, z)$
的全部分析流程，是理解非连通胶子三点函数分析的优秀范例。

\subsection{代码物理框架速查表}

\begin{table}[H]
\centering
\caption{物理公式 $\leftrightarrow$ huangcl 代码行号速查}
\begin{tabular}{c|c|c}
\toprule
\textbf{物理内容} & \textbf{公式编号} & \textbf{代码行号} \\
\midrule
OPE 非极化组合 & (\ref{eq:O_euclidean_components}) & 89 \\
时间平移 (绝对 $\to$ 相对) & -- & 92--104 \\
非连通 3pt 构造 & (\ref{eq:C3pt_disc}) & 106--114 \\
Jackknife 重采样 & (\ref{eq:jackknife}) & 53--64, 120--122 \\
非连通减除 & (\ref{eq:C3pt_disc}) & 127--129 \\
比值 $R = C_3/C_2$ & (\ref{eq:ratio_def}) & 130--132 \\
Jackknife 误差 & (\ref{eq:jackknife_error}) & 46--50 \\
比值拟合 (两态模型) & (\ref{eq:ratio_fit}) & (后续步骤) \\
\bottomrule
\end{tabular}
\end{table}


% ============================================================
% 附录：数据路径
% ============================================================
\appendix
\section{数据路径与参考源}

\subsection{集群数据路径}

\begin{table}[H]
\centering
\caption{集群上的数据路径}
\begin{tabular}{c|l}
\toprule
\textbf{数据类型} & \textbf{路径模式} \\
\midrule
2pt 关联函数 & \texttt{/public/group/lqcd/donghx/2pt\_Result/} \\
 & \quad \texttt{\{conf\_name\}/momsmear2z/\{conf\_id\}/} \\
 & \quad \texttt{twopt\_slice\_pp\_...\_nopol\_ss\_conf\{id\}.npy} \\
\midrule
OPE 算符 & \texttt{/public/group/lqcd/donghx/Ope\_Gluon/} \\
 & \quad \texttt{Result\_hpy\_4D\_10times/\{conf\_short\}/zdir/} \\
 & \quad \texttt{\{conf\_id\}/ops\_mu\{mu\}\_nu\{nu\}\_dz24\_conf\{id\}.npz} \\
\midrule
输出 & \texttt{/public/home/zhangxin/lattice-pdf/} \\
 & \quad \texttt{examples/huangcl/result/ratio.png} \\
\bottomrule
\end{tabular}
\end{table}

\subsection{参考文献}

\begin{thebibliography}{99}

\bibitem{gluon_pdf_derivation}
张鑫. \textit{胶子 PDF 的理论推导}.
\texttt{/root/lattice-pdf/文档/gluon\_pdf\_derivation.tex}.

\bibitem{note_of_gluon_PDFs}
张鑫. \textit{Note of Gluon PDFs}.
\texttt{/root/lattice-pdf/文档/note\_of\_gluon\_PDFs.tex}.

\bibitem{场强张量}
张鑫. \textit{格点 QCD 中的场强张量}.
\texttt{/root/lattice-pdf/补充/格点QCD中的场强张量.tex}.

\bibitem{关联函数}
张鑫. \textit{格点 QCD 中的关联函数}.
\texttt{/root/lattice-pdf/补充/格点QCD中的关联函数.tex}.

\bibitem{大动量有效理论}
张鑫. \textit{格点 QCD 中的大动量有效理论 (LaMET)}.
\texttt{/root/lattice-pdf/补充/格点QCD中的大动量有效理论.tex}.

\bibitem{蒸馏方法}
张鑫. \textit{格点 QCD 蒸馏方法解析}.
\texttt{/root/lattice-pdf/补充/格点QCD蒸馏方法解析.tex}.

\bibitem{光锥PDF}
张鑫. \textit{格点 QCD 中的光锥 PDF 与 quasi-PDF}.
\texttt{/root/lattice-pdf/补充/格点QCD中的光锥PDF与quasi\_PDF.tex}.

\bibitem{lamet_original}
X. Ji, \textit{Parton Physics on a Euclidean Lattice},
Phys. Rev. Lett. 110, 262002 (2013).

\bibitem{distillation}
M. Peardon et al. (Hadron Spectrum Collaboration),
\textit{A Novel Quark-Field Creation Operator Construction for Hadronic Physics in Lattice QCD},
Phys. Rev. D 80, 054506 (2009).

\bibitem{gluon_pdf_first}
Z. Fan, Y. Yang, et al.,
\textit{First Nucleon Gluon PDF from Large Momentum Effective Theory},
Phys. Rev. Lett. 127, 062002 (2021).

\bibitem{gluon_pdf_continuum}
Z. Fan et al.,
\textit{First Self-Renormalized Gluon PDF of Nucleon from Large-Momentum Effective Theory in the Continuum Limit},
arXiv:2401.xxxx (2024).

\bibitem{hybrid_renorm}
X. Ji, Y. Liu, et al.,
\textit{A Hybrid Renormalization Scheme for Quasi Light-Front Correlations in Large-Momentum Effective Theory},
Nucl. Phys. B 964, 115311 (2021).

\bibitem{claude_md}
\texttt{/root/lattice-pdf/CLAUDE.md} —— 项目总览与代码约定.

\end{thebibliography}

\subsection{相关代码文件清单}

\begin{table}[H]
\centering
\caption{本文引用的源代码文件}
\begin{tabular}{c|l}
\toprule
\textbf{文件} & \textbf{功能} \\
\midrule
\texttt{examples/huangcl/code.py} & 本文分析的主体代码 \\
\texttt{examples/huangcl/submit.sh} & Slurm 作业提交脚本 \\
\texttt{examples/donghx/Calc\_ope\_unpol.py} & OPE 算符的 MPI 并行计算 \\
\texttt{examples/donghx/Operator.py} & Plaquette, Clover, OPE 算符核心函数 \\
\texttt{examples/donghx/2pt\_proton\_Cg5gmu\_L24x72\_mom2\_zdir\_gpu.py} & 蒸馏法核子 2pt 计算 \\
\texttt{examples/donghx/gamma\_matrix\_cupy\_DR.py} & DeGrand-Rossi 基 gamma 矩阵 \\
\bottomrule
\end{tabular}
\end{table}

\end{document}
